\uuid{k6SR}
\exo7id{827}
\titre{exo7 827}
\auteur{ridde}
\organisation{exo7}
\datecreate{1999-11-01}
\isIndication{false}
\isCorrection{false}
\chapitre{Calcul d'intégrales}
\sousChapitre{Fraction rationnelle}
\cpge{mpsi-mp2i.polynomes-fractions-rationnelles.i, pcsi.polynomes.f, ptsi.polynomes.f, mpsi-mp2i.primitives-equations-differentielles.a, pcsi.techniques-calcul-primitives.a, ptsi.techniques-calcul-primitives.a, bcpst1.derivees-primitives-integrales.c}
\module{Analyse}
\niveau{L1}
\difficulte{}

\contenu{
\question{
D\'eterminer les intervalles d'\'etude et calculer les primitives des fonctions :
$$ \frac 1{ (x + 2) (x^2 + 2x + 5)}$$
$$ \frac{2x}{ (1-x + x^2)^2}$$
$$ \frac{x^2}{ (x-1)^2 (x^2 + 4)}$$
$$ \frac 1{ (1 + x^3)^3}$$
}
}
